One of the fundamental problems in reinforcement learning is credit assignment, that is, determining the contribution of individual actions to the final reward. This problem is particularly pronounced in tasks with terminal rewards, where the entire episode is evaluated by a single scalar signal and no intermediate feedback is available. In such settings, the algorithm receives no information about which actions contributed positively or negatively to the outcome, leading to high gradient variance and making learning difficult, especially for long action sequences.

The recently proposed method SCAR (Shapley Credit Assignment for More Efficient RLHF) formulates sequence generation as a cooperative game and uses Shapley values to decompose a terminal reward into contributions of individual tokens. In the context of RLHF, this makes it possible to replace a single sparse reward with a dense sequence of token-level rewards, which can be used for more stable and efficient fine-tuning of large language models.

In this thesis, the idea of Shapley-based credit assignment is transferred from RLHF to the Dreamer algorithm. Unlike RLHF, where reward evaluation requires a separate reward model trained on human preferences, Dreamer already includes a learned model of environment dynamics and reward. This allows counterfactual variants of trajectories to be evaluated using latent imagined rollouts, without additional interaction with the real environment.

The work focuses on text-based and text--procedural tasks in which an agent executes a sequence of discrete actions (such as text generation, memory operations, or module calls) and receives a reward only at the end of the episode. After an episode is completed, the terminal reward is distributed among individual actions using Shapley sampling. This is achieved by estimating the expected terminal reward for multiple modified variants of the same trajectory, obtained by replacing a subset of actions with a baseline action, and evaluating them through latent rollouts of the Dreamer environment model. The resulting sequence of action contributions is used as a dense reward signal for training both the policy and the dynamics model.

In the experimental part, Dreamer is compared under three training regimes: using only a terminal reward, using manual reward shaping, and using Shapley-based reward decomposition. The evaluation focuses on learning speed, training stability, and the quality of the learned policy, as well as on the impact of the number of Shapley samples and the action segmentation strategy.